\section{The exact logarithmic scale of the finite normalization}
\label{sec:spectral-endpoint-asymptotic}

\begin{theorem}[AC-S41: a nonzero finite endpoint asymptotic]
\label{AC-S41}
Keep the fixed-window hypotheses and phase convention of AC-S40.
There is a real constant \(\kappa_\lambda\ne0\) such that
\[
             \sqrt{\log N}\,u_N(\ell)
                    \longrightarrow\kappa_\lambda,
                 \qquad \ell=2\log\lambda.               \tag{FA.1}
\]
It has the following exact description.  Let
\(R_o=(A_{\rm odd}-\alpha)^{-1}\), and choose any real
\(g\in C_c^\infty(0,\ell)\), odd about \(\ell/2\), for which
\(\tau_+(R_og)\ne0\).  Such tests exist.  Then
\[
 \kappa_\lambda
       =\frac{\displaystyle\int_0^\ell u(x)g'(x)dx}
                    {2\tau_+(R_og)}.                    \tag{FA.2}
\]
The ratio is independent of the chosen test.
In particular, for the original source normalization
\(\xi_N=u_N/u_N(\ell)\),
\[
 \frac{\xi_N}{\sqrt{\log N}}
           \longrightarrow\frac{u}{\kappa_\lambda}
                  \quad\hbox{in }L^2\hbox{ and form norm},
                                                        \tag{FA.3}
\]
and the actual S25 determinants satisfy
\[
 \frac{\det_{\rm reg}(D^{(\lambda,N)}_{\log}-z)}
                           {\sqrt{\log N}}
      \longrightarrow
       -i\lambda^{-iz}\frac{\widehat u(z)}{\kappa_\lambda}
                 \quad\hbox{locally uniformly in }\mathbb C.
                                                        \tag{FA.4}
\]
The coefficient in (FA.2) is specified by the actual window
operator.  AC-S43 below identifies it exactly with the renewal trace
\(\tau_+(u)\).  The theorem is unconditional on
\(1<\lambda\leq\sqrt{r_0}\), \(r_0^3-r_0-1=0\).
\end{theorem}
\begin{proof}
Use the inner product linear in its second argument.
For real odd \(g\in C_c^\infty(I)\), put
\[
 w=R_og,\qquad \phi=iw,\qquad f_0=ig,
 \qquad D=\frac{\ell}{2\pi i}\frac{d}{dx},
 \qquad Df_0=\frac{\ell}{2\pi}g'.                         \tag{FA.5}
\]
The bounded odd-resolvent argument in step~1 of AC-S40 places
\(w,\phi\) in \(\mathcal X_I\); in particular their traces exist.
The same step constructs at least one such \(g\) with
\(\tau:=\tau_+(w)\ne0\).
The Fourier coefficients of \(\phi\) and \(f_0\) are real and
odd, while those of \(u_N\) are real and even.

Write \(c_N=\sqrt\ell\,u_N(\ell)\).
AC-S40 proves that \(c_N\sqrt{\log N}\) is bounded.
For the fixed test with \(\tau\ne0\), the exact displacement
pairing (FE.20), (FE.26) is
\[
 -c_N\langle\beta_N,\phi\rangle
   =\frac{\ell}{2\pi}\int_0^\ell u_N(x)g'(x)dx+\varepsilon_N,
                    \qquad \varepsilon_N=o(1).            \tag{FA.6}
\]
The error is indeed \(o(1)\), because (FE.27) is
\(o(|c_N|\sqrt{\log N})\) and this last factor is bounded.
There is no conjugation ambiguity in (FA.6): all the relevant
Fourier coefficients in this pairing are real, and
\(\langle P_N f_0,D_Nu_N\rangle
 =\langle Df_0,u_N\rangle\).
Using (FE.23) and \(u_N\to u\), equation (FA.6) gives a limit
\[
 \lim_{N\to\infty}c_N\sqrt{\log N}
   =\frac{\sqrt\ell}{2\tau}\int_0^\ell u(x)g'(x)dx.
                                                        \tag{FA.7}
\]
This proves existence of (FA.1), with the value (FA.2),
but has not yet shown that this value is nonzero.

For that purpose take any real odd compact smooth \(g\), allowing
\(\tau_+(R_og)=0\).  The same construction and Fourier-tail
estimates apply to \(\phi=iR_og\).
The Abelian estimate in the proof of (FE.23) also applies when
the trace is zero: the boundary function is then
\(o((\log(1/d))^{-1/2})\), and the stated remainder estimate
gives
\[
 \frac{\langle\beta_N,\phi\rangle}{\sqrt{\log N}}
           \longrightarrow-\frac{\sqrt\ell}{\pi}
                                         \tau_+(R_og).
                                                        \tag{FA.8}
\]
Consequently (FA.6), now with this arbitrary test, implies
the exact identity
\[
              \int_0^\ell u(x)g'(x)dx
                        =2\kappa_\lambda\tau_+(R_og).
                                                        \tag{FA.9}
\]
If \(\kappa_\lambda=0\), the left side vanishes for every
real odd compact smooth \(g\).
For every real even compact smooth test \(h\) of integral zero,
the function \(g(x)=\int_0^x h(y)dy\) is compactly supported,
smooth, and odd about \(\ell/2\).  Thus \(\int uh=0\).
Fix an even compact smooth test of integral one and subtract its
appropriate multiple from any other even test.
It follows that the distribution defined by \(u\) is constant
on even tests.  It vanishes on odd tests by evenness of \(u\),
so \(u\) is a constant distribution on the whole interval.
The continuous representative of AC-S36 has zero endpoint values;
therefore this constant is zero, contrary to \(\|u\|_2=1\).
We have proved \(\kappa_\lambda\ne0\).
Equation (FA.9) also proves the independence in (FA.2).

Finally, divide the already proved \(L^2\) and form-norm
convergence \(u_N\to u\) by the nonzero scalars
\(\sqrt{\log N}\,u_N(\ell)\to\kappa_\lambda\).
This proves (FA.3).
Compact support and Cauchy--Schwarz imply locally uniform entire
convergence of the Fourier transforms: on every compact subset
of \(\mathbb C\), the Fourier kernels have uniformly bounded
\(L^2(I)\) norms.
Apply the literal source identity
\(\det_{\rm reg}(D^{(\lambda,N)}_{\log}-z)
 =-i\lambda^{-iz}\widehat\xi_N(z)\)
to obtain (FA.4).
\end{proof}

\paragraph{What has and has not been normalized.}
The factor \(\sqrt{\log N}\) in (FA.4) concerns the original
unsmoothed finite evaluation, with the source's fixed Fourier
cutoff \(|n|\leq N\).  It is now proved, rather than imposed by
the semigroup construction of AC-S39.  Formula (FA.2) fixes its
nonzero coefficient through the actual resolvent.
AC-S43 identifies this coefficient exactly with the nonzero continuum
renewal trace.  AC-S42 extends the finite asymptotic to every fixed
simple-even window.  Proving that
lowest-line hypothesis for all windows and the separate large-window
ground-line selection problem remain open.
